% =====================================================================================================
% Section III. The Technology from Designed Variation (sections/technology.tex), version 2.
% Sources: output/memos/ra1_modelfree.md (+_review), m2_techpanel.md (+_review), m8_measurement.md (+_review),
% ra3_econ.md; output/tables/ra1_modelfree_*.csv, m2_*.csv, ra3_econ_data_demand.csv; paper/notes/m9_preanalysis_plan.md;
% code/sweep/run_grid.py. Number register, referee comments and open issues: paper/notes/writer2_technology.md.
% =====================================================================================================

\section{The Technology from Designed Variation}\label{sec:tech}

Section~\ref{sec:ident} showed that IsoFLOP profiles, which hold compute fixed and vary the input mix, identify the
curvature of the technology without a functional form (Proposition~\ref{prop:modelfree}). This section takes that
result to the data and establishes three facts. First, on every public IsoFLOP design with annealed runs, the
elasticity of substitution on the expansion path, $\sigma^*$, lies between 0.66 and 0.71, and the estimates from three
developers' independent designs are consistent with a common value of about 0.70. Second, parametric scaling laws
match this benchmark, where the data can tell them apart, only when the outer exponent $\kappa$ is free. Third, the
other objects that the inversion of Section~\ref{sec:wedge} needs, the allocation exponent $a$ and the compute-optimal
ratio $M^*$, are not portable across designs, conventions and data. Subsection~\ref{sec:tech:experiment} describes
our own experiment and its pre-registered analysis.

\subsection{Designs}\label{sec:tech:designs}\label{sec:tech:chinchilla}

Table~\ref{tab:designs} lists the designs. The most informative are IsoFLOP designs: at each of several compute
budgets, the experimenter trains models of different sizes on the isocost line $6ND=C$. Each budget yields a profile
of loss against $\ln N$ whose minimum lies on the expansion path; these minima are outcomes of assigned inputs, not
choices.

Four studies provide such profiles. Epoch AI's digitization of the runs in Figure~4 of \citet{hoffmann2022training},
the sample of \citet{besiroglu2024chinchilla}, has 245 runs with training compute from $1.4\times10^{18}$ to
$1.3\times10^{22}$ FLOP: 137 on nine IsoFLOP profiles from $6\times10^{18}$ to $3\times10^{21}$ FLOP, and 108 from
35 fixed-size trunks with different training horizons. Model-free estimates use the 137 profile runs; parametric fits
use the full design without the five highest-loss runs ($n=240$). Meta's Llama~3 profiles \citep{grattafiori2024llama}, digitized by \citet{czech2026llama3isoflop}, cover ten
budgets up to $10^{22}$ FLOP; their unstated loss units do not matter, because the model-free estimator is invariant
to monotone transformations of output. Marin's released ladders \citep{marin2026ladders} repeat the design on three
corpora with budgets up to $3\times10^{20}$ FLOP. \citet{porian2024resolving} release loss curves of 16 architectures
trained at a constant learning rate, from which we read IsoFLOP-like profiles at 12 budgets up to
$2.6\times10^{19}$ FLOP; because the runs are not annealed, these are losses of unfinished schedules.

Farseer \citep{li2025predictableb} is not an IsoFLOP design, but its 404 released runs (of about 1,000 trained) form
a near-factorial grid dense enough to locate the expansion path inside it. Three smaller sweeps
of 30--35 runs have no IsoFLOP profiles: the three web corpora of \citet{gadre2024language}, all evaluated on C4, the
OLMo ladder \citep{bhagia2024establishing}, and the single-epoch runs of \citet{muennighoff2023scaling}. With
Chinchilla's full design and Farseer, they form the seven sweep--corpus technologies from five studies on which we
compare functional forms. DataDecide \citep{magnusson2025datadecide}, whose final checkpoints all have $D\approx100N$,
serves only to compare data recipes.

The model-free estimator requires each profile to be an exact isocost, so $N$ is measured in the convention that
makes $C=6ND$ hold along it (Table~\ref{tab:designs}, notes), and only unit-free objects are compared across rows. The
off-path spread, the residual standard deviation of $\ln M$ given $\ln C$ with $M=D/N$, is zero for compute-optimal
runs and by Proposition~\ref{prop:ident} governs how much a design reveals about curvature: 1.4--2.5 in the IsoFLOP
designs and 0.94 on the OLMo ladder.

% Table 2 (Section III): designs and the elasticity of substitution on the expansion path.
% Sources (every cell): Panel A model-free and same-run parametric columns: output/tables/ra1_modelfree_isoflop_summary.csv
% (sigma_re, se_sigma_re), ra1_modelfree_isoflop_param.csv (sigma_kappa, sigma_chin, FHH wild SEs); Farseer model-free:
% ra1_modelfree_farseer_sigma_slope.csv (first-derivative estimate, 0.7078; 0.7008-0.7114 across bandwidths); Farseer and
% Panel B parametric: m2_table3_technology.csv (Huber rows; sigma_star, sigma_star_q with pairs/cell bootstrap SEs).
% Runs, N, M, off-path spread: recomputed from the ra1 loaders (ra1_common.isoflop_designs, sweep_panels); sweep spreads
% equal m2's sd_offpath. Panel C: ra1_modelfree_heterogeneity.csv (rows "all technologies" of the named sets).
% Writer notes: paper/notes/writer2_technology.md.
\begin{table}[tp]
\centering
\caption{Designs and the Elasticity of Substitution on the Expansion Path}
\label{tab:designs}
\footnotesize
\setlength{\tabcolsep}{3pt}
\begin{tabular*}{\textwidth}{@{\extracolsep{\fill}}lrcccccc@{}}
\toprule
 & & & & Off-path & \multicolumn{3}{c}{$\sigma^*$} \\
\cmidrule(l){6-8}
Design & Runs & $N$ & $M=D/N$ & spread & Model-free & $\kappa$ free & $\kappa=1$ \\
\midrule
\multicolumn{8}{@{}l}{\textit{Panel A. IsoFLOP designs and Farseer's factorial grid}}\\[1pt]
Chinchilla, IsoFLOP runs & 137 & 57M--16B & 0.04--294 & 1.67 & 0.673 & 0.667 & 0.692 \\
 & & & & & (0.027) & (0.023) & (0.013) \\
Llama 3 & 133 & 59M--17B & 0.44--643 & 1.38 & 0.660 & 0.693 & 0.769 \\
 & & & & & (0.023) & (0.020) & (0.006) \\
Marin, Comma & 85 & 0.12--12.6B & 0.09--3,706 & 2.46 & 0.700 & 0.678 & 0.690 \\
 & & & & & (0.027) & (0.015) & (0.015) \\
Marin, DCLM & 85 & 0.12--12.6B & 0.09--3,706 & 2.46 & 0.713 & 0.663 & 0.673 \\
 & & & & & (0.027) & (0.013) & (0.013) \\
Marin, Nemotron-CC & 88 & 0.12--12.6B & 0.09--3,706 & 2.42 & 0.705 & 0.677 & 0.684 \\
 & & & & & (0.023) & (0.014) & (0.014) \\
Porian, RefinedWeb$^{a}$ & 121 & 5.2M--0.90B & 0.02--346 & 2.29 & 0.518 & 0.345 & 0.599 \\
 & & & & & (0.017) & (0.141) & (0.032) \\
Porian, OpenWebText2$^{a}$ & 116 & 5.2M--0.90B & 0.02--194 & 2.20 & 0.505 & 0.347 & 0.609 \\
 & & & & & (0.014) & (0.160) & (0.039) \\
Farseer & 404 & 0.10--6.4B & 0.31--2,570 & 1.97 & 0.708$^{b}$ & 0.710 & 0.772 \\
 & & & & &  & (0.003) & (0.010) \\[3pt]
\multicolumn{8}{@{}l}{\textit{Panel B. Other sweep--corpus technologies (parametric fits to all runs)}}\\[1pt]
Chinchilla, all runs & 240 & 57M--16B & 0.46--341 & 1.35 & -- & 0.701 & 0.737 \\
 & & & & &  & (0.014) & (0.006) \\
Gadre, C4 & 34 & 11M--6.9B & 5--640 & 1.46 & -- & 0.607 & 0.795 \\
 & & & & &  & (0.054) & (0.020) \\
Gadre, RedPajama & 35 & 11M--6.9B & 5--640 & 1.47 & -- & 0.623 & 0.795 \\
 & & & & &  & (0.059) & (0.013) \\
Gadre, RefinedWeb & 35 & 11M--6.9B & 5--640 & 1.47 & -- & 0.594 & 0.796 \\
 & & & & &  & (0.062) & (0.015) \\
OLMo ladder & 30 & 0.19--3.2B & 10--200 & 0.94 & -- & 0.544 & 0.825 \\
 & & & & &  & (0.021) & (0.034) \\
Muennighoff & 33 & 7.1M--8.7B & 0.04--303 & 1.97 & -- & 0.511 & 0.735 \\
 & & & & &  & (0.056) & (0.031) \\
\bottomrule
\end{tabular*}

\vspace{4pt}
\begin{tabular*}{\textwidth}{@{\extracolsep{\fill}}lccccc@{}}
\multicolumn{6}{@{}l}{\textit{Panel C. Heterogeneity across designs (random effects)}}\\[1pt]
\midrule
Set & $k$ & Mean [95\% CI] & Between s.d. & $I^2$ & Cochran $Q$ ($p$) \\
\midrule
Model-free, Panel A excluding Porian & 6 & 0.695 [0.673, 0.717] & 0.000 & 0.00 & 4.1 (0.54) \\
Model-free, all of Panel A & 8 & 0.646 [0.562, 0.730] & 0.098 & 0.96 & 188.9 ($<$0.001) \\
$\kappa$ free, IsoFLOP designs of Panel A & 7 & 0.672 [0.636, 0.707] & 0.019 & 0.47 & 11.4 (0.08) \\
$\kappa$ free, sweep--corpus technologies & 7 & 0.622 [0.549, 0.695] & 0.068 & 0.93 & 82.1 ($<$0.001) \\
$\kappa=1$, sweep--corpus technologies & 7 & 0.777 [0.745, 0.809] & 0.030 & 0.83 & 36.3 ($<$0.001) \\
\bottomrule
\end{tabular*}

\begin{tablenotes}
$N$: parameters in the convention that makes each IsoFLOP profile an exact isocost (Chinchilla: total; Llama~3,
Marin: $C/(6D)$; Porian et al.: non-embedding plus output head); Farseer non-embedding; OLMo ladder without the
input embedding; Gadre et al.\ total; Muennighoff et al.\ the authors' count. Output: validation loss on MassiveText
(Chinchilla), in unstated units (Llama~3), on Paloma (Marin), on the training corpus (Porian et al.), in bits per
character (Farseer), on C4 (others). Off-path spread: residual s.d.\ of $\ln M$ given $\ln C$. Model-free: equation
\eqref{eq:ident-mf} per budget, random-effects mean over budgets, design-conditional wild-bootstrap standard errors.
Parametric columns: Huber loss on $\ln L$; Panel A on the same runs (Feng--He--Hu wild standard errors), Farseer and
Panel B pairs or cell bootstrap. Panel C: DerSimonian--Laird mean with Hartung--Knapp--Sidik--Jonkman interval;
sweep--corpus = Farseer and Panel B.
$^{a}$Losses read along constant-learning-rate runs without cooldown.
$^{b}$First-derivative estimate on the local expansion path (0.701--0.711 across bandwidths); Panel C and
Figure~\ref{fig:sigma} use the Hessian-based estimate, 0.703 (0.013).
\end{tablenotes}

\begin{tablenotes}[Source]
\citet{hoffmann2022training}; \citet{besiroglu2024chinchilla}; \citet{grattafiori2024llama};
\citet{czech2026llama3isoflop}; \citet{marin2026ladders}; \citet{porian2024resolving}; \citet{li2025predictableb};
\citet{gadre2024language}; \citet{bhagia2024establishing}; \citet{muennighoff2023scaling}; authors' calculations
(Online Appendix~\ref{app:additional}).
\end{tablenotes}
\end{table}


\subsection{The Elasticity without a Functional Form}\label{sec:tech:modelfree}

By Proposition~\ref{prop:modelfree}(ii), $1/\sigma^*-1$ equals the curvature of the IsoFLOP profile in $n=\ln N$ at
its minimum divided by twice the slope of the IsoFLOP-minima frontier, the curve traced by the minima
(equation~\ref{eq:ident-mf}). Both are in the same units of output, so the irreducible loss $E$, the outer exponent and
the loss units drop out.

We estimate both terms budget by budget. A quadratic in $\ln N$ fitted to the runs within one log point of the path
gives the curvature and the minimum, a cubic in $\ln C$ through the logged minima gives the frontier slope, and a
design's estimate is the random-effects mean over budgets \citep{dersimonian1986meta}. Standard errors come from a
design-conditional wild bootstrap that holds the design fixed and re-estimates every step in each of 999 draws. On
Farseer we fit a kernel-weighted local quadratic of $\ln L$ in $(\ln N,\ln D)$, find the points where the local wedge
$w$ equals one on eight isocosts, and read $\sigma^*$ off the slope of $\ln w$ in $\ln M$ there, which equals
$1/\sigma^*-1$ for any smooth technology (Proposition~\ref{prop:modelfree}(iii)). Online
Appendix~\ref{app:additional} gives details and robustness.

\begin{figure}[tp]
\centering
\includegraphics[width=\textwidth]{ra1_modelfree_sigma_by_design}
\caption{The Elasticity of Substitution by Design, with and without a Functional Form}\label{fig:sigma}
\begin{figurenotes}
Panel a: $\sigma^*$ by design with 95 percent intervals ($\pm1.96$ standard errors). Circles: model-free estimate,
equation~\eqref{eq:ident-mf}, random-effects mean over budgets; diamonds and squares: $\kappa$ family and Chinchilla
form ($\kappa=1$) fitted by Huber loss to the same runs. Farseer: Hessian-based estimate on the local expansion path
(0.703; the first-derivative estimate is 0.708) and full-grid parametric fits. Panel b: model-free estimate by budget
with 95 percent bootstrap intervals; solid line, design mean; dashed, $\kappa=1$; dotted, $\kappa$ free. Farseer
(last panel): local-path estimates at eight compute levels, with the path value implied by Farseer's own functional
form (gray). Porian et al.'s losses come from runs without cooldown.
\end{figurenotes}
\begin{figurenotes}[Source]
Designs of Table~\ref{tab:designs}; authors' calculations.
\end{figurenotes}
\end{figure}

Figure~\ref{fig:sigma} and Table~\ref{tab:designs} report the results. On Chinchilla's profiles $\sigma^*=0.673$
(standard error 0.027), on Meta's 0.660 (0.023), and on Marin's three corpora 0.700, 0.713 and 0.705. Farseer's local
expansion path gives 0.708, and 0.701--0.711 at every bandwidth we try. The six estimates are consistent with a common
value. Their random-effects mean is 0.695 (95 percent confidence interval [0.673, 0.717]), the estimated
between-design standard deviation is zero, and Cochran's homogeneity statistic is 4.1 on five degrees of freedom
($p=0.54$). With one estimate per study the mean is 0.693. The test has little power against a between-design
standard deviation below about 0.02, so ``homogeneous'' means within two or three hundredths of 0.70. The summary
enters Farseer's Hessian-based estimate, 0.703, which is sensitive to the bandwidth: at the narrower cross-validated
bandwidth it is 0.664, and the mean 0.679.

The estimator is accurate on these grids: on each design's own fitted technology it is biased by at most 0.006, and its
interval covers the truth in 92--98 percent of Monte Carlo replications (Online Appendix~\ref{app:add-modelfree}). The
larger threat is FLOP accounting: if FLOPs per token of the measured parameter count vary with size with elasticity
$\eta$, $\sigma^*$ shifts by about $0.44\eta$. In Marin's configuration parameter count ($\eta=0.087$) its estimates fall
to 0.64--0.66; Chinchilla's and Meta's FLOP accounting is unobserved, which adds uncertainty of about $\pm0.04$ for
$|\eta|\le0.1$.

Porian et al.'s profiles are the exception, at 0.518 (0.017) and 0.505 (0.014), and we keep them out of the summary.
Their losses come from constant-learning-rate runs without cooldown, and the gap between unannealed and annealed loss
varies with $D/N$ at fixed compute, which bends the profile: the estimates move between 0.45 and 0.55 with the window.
They also come from the smallest models and budgets in the table. Including them lowers the mean to 0.646 and makes
the set strongly heterogeneous (Table~\ref{tab:designs}, panel C). Our experiment, with annealed runs at a similar
scale, can separate schedule from scale (subsection~\ref{sec:tech:experiment}).

Within designs, $\sigma^*$ falls with compute: by 0.072 per decade in Chinchilla (standard error 0.026), by 0.053 in
Llama~3 (0.010) and by 0.015 in Farseer (0.008), but not detectably in Marin. Neither parametric family allows such
drift, and the estimator shows none when the truth has none, so this is evidence of non-homotheticity, although weak:
across windows the Chinchilla drift is 0.034--0.096 per decade. If it is real, $\sigma^*$ at frontier compute is
below 0.70.

That $\sigma^*<1$ is not a finding: by equation~\eqref{eq:ident-mf} it holds exactly when IsoFLOP profiles are U-shaped
under a falling frontier (Lemma~\ref{lemma:sigma}). The level matters because it prices departures from the
compute-optimal mix. With $\sigma^*=0.70$ (and equal
exponents), a model trained on ten times the compute-optimal tokens per parameter needs 73 percent more compute than
the compute-optimal model of the same loss, against 57 percent at $\sigma^*=0.74$ (Online
Appendix~\ref{app:additional}). A developer who makes that choice reveals a value of compactness 1.7 times its
training expenditure ($w=2.7$, $s=0.63$; Proposition~\ref{prop:wedge}). Because $D$ counts tokens processed,
$\sigma^*$ is an elasticity between parameters and training tokens, not between compute and a stock of unique text
(Section~\ref{sec:econ}).

\subsection{Parametric Forms}\label{sec:tech:parametric}\label{sec:tech:sweeps}

Scaling laws are usually fitted in the Chinchilla form, the $\kappa=1$ member of the family
$L=E+[AN^{-\alpha}+BD^{-\beta}]^{\kappa}$ of Section~\ref{sec:framework}. In every member $\sigma^*=2/(2+\alpha+\beta)$
depends only on the inner exponents, but imposing $\kappa=1$ ties the curvature of IsoFLOP profiles to the slopes of
the frontier and the path (Proposition~\ref{prop:modelfree}(iv)). The restriction is rejected wherever it can be
tested. On Chinchilla's full design $\hat\kappa=0.774$, and $\kappa=1$ is rejected at 1 percent under five of the
six inference schemes we consider, including a cluster bootstrap over the nine budgets and the 35 off-profile trunks
($p=0.004$); the exception, the original clustering that assigns off-profile runs to the nearest budget, gives
$p=0.052$. It is rejected in
all seven sweep--corpus technologies ($p<0.001$) and on every IsoFLOP design but Chinchilla's profile runs, where the
quasi-likelihood-ratio statistic of 7.8 is borderline.

Rejection concerns fit; the question here is whether the restriction distorts $\sigma^*$. The model-free estimates
answer it on the same runs (Figure~\ref{fig:sigma}, panel a). Where the comparison discriminates, the $\kappa$-free
form agrees with the benchmark and the Chinchilla form does not. On Llama~3, $\kappa=1$ gives 0.769, 4.8 standard
errors above the model-free 0.660, while $\kappa$ free gives 0.693, 1.5 standard errors away. On Farseer, $\kappa=1$
gives 0.772, against 0.708 on the local path and 0.710 with $\kappa$ free. Elsewhere the comparison does not
discriminate: on Chinchilla's profile runs both forms lie within one standard error of the benchmark, on Marin the
Chinchilla form is slightly closer, and on Porian et al.'s profiles the benchmark lies between the two. On
Chinchilla's full design the two forms give 0.737 and 0.701, and the bootstrap-calibrated profile-likelihood set for
$\sigma^*$ with $\kappa$ free, [0.68, 0.72], excludes 0.74 ($p=0.003$, the floor with 299 draws); the contrast between
0.737 and the model-free 0.673 compares different samples. The form's other restriction, equal exponents, survives
at 5 percent in every sweep under tests robust to heavy-tailed residuals (Online Appendix~\ref{app:additional}).

With $\kappa$ free, the three sweeps without IsoFLOP profiles give lower values: 0.59--0.62 on Gadre et al.'s
corpora, 0.544 on the OLMo ladder and 0.511 on Muennighoff et al.'s runs. Across the seven sweep--corpus technologies
the estimates are strongly heterogeneous: Cochran's statistic is 82.1 on six degrees of freedom ($p<0.001$), the
between-technology standard deviation is 0.068, and the random-effects mean is 0.622 (95 percent confidence interval
[0.549, 0.695]). The heterogeneity survives dropping the OLMo ladder or keeping one estimate per study. Chinchilla and
Farseer agree ($p=0.51$), and each differs from the OLMo ladder and from Muennighoff et al.\ by 3.3 to 7.8 standard
errors. On the seven IsoFLOP designs, by contrast, the $\kappa$-free fits are close to homogeneous ($p=0.08$).

We cannot tell whether the low values reflect the technology or the $\kappa$ family's fit to small designs dominated
by off-path runs. With 30--35 runs, $E$ and $\kappa$ are hard to separate: over the confidence set for $E$,
$\sigma^*$ moves by up to 0.08 on these sweeps, against at most 0.006 on Chinchilla, Farseer and the OLMo ladder
(Online Appendix Table~\ref{tab:app-kappa-robust}). The OLMo ladder is precise and still low, also with task bits per
byte as output (0.51), so a difference in recipe cannot be ruled out. The Chinchilla form raises every estimate, to
0.735--0.825. One more
restriction of the family fails where it can be checked: the log wedge is linear in $\ln(M/M^*(C))$ in every member,
but convex inside Farseer's grid (Section~\ref{sec:wedge}).

In sum, where IsoFLOP profiles from annealed runs exist, $\sigma^*$ is about 0.70 with or without a functional form;
small sweeps fitted parametrically give 0.51--0.62; and the Chinchilla restriction overstates $\sigma^*$ wherever the
data can detect it. Section~\ref{sec:wedge} therefore uses the $\kappa$-free Chinchilla technology as its reference.

\subsection{What Is Not Portable}\label{sec:tech:portable}\label{sec:tech:quality}\label{sec:tech:measurement}

The inversion of Section~\ref{sec:wedge} needs a scale, $1/\sigma^*-1$, and a zero point, $M^*(C)$. The scale is
estimated above. The zero point, and the allocation exponent $a$ that carries it across compute, are where designs
disagree. Under the Chinchilla form, $a$ ranges from 0.37 (OLMo ladder) to 0.57 (Gadre et al., RefinedWeb) across the
seven sweep--corpus technologies, and from 0.36 to 0.57 once lab laws and Farseer's own functional form are added
(Section~\ref{sec:econ}). The compute-optimal ratio $M^*(10^{21})$ ranges from 3.4 to 60 tokens per parameter, an
18-fold spread. Even the model-free minima paths differ, with slopes of 0.50 on Chinchilla and Llama~3 and 0.35--0.43
on Marin's three corpora.

$M^*$ is a level: it depends on the ratio $A/B$, which absorbs every difference in tokenizers, parameter counts and
recipes, and $a$ is its slope in compute. Counting Farseer's embedding parameters moves $a$ from 0.53 to 0.41 and
$M^*(10^{23})$ from 19.5 to 45; counting non-embedding parameters on Chinchilla raises $a$ from 0.514 to 0.556. This is
not evidence that $\sigma^*$ is intrinsically better determined than $a$: the same change of convention moves
Farseer's $\kappa$-free $\sigma^*$ from 0.710 to 0.668. The Kaplan--Chinchilla disagreement is the canonical case. On
the 975 released runs of \citet{porian2024resolving}, the allocation exponent falls from 0.84 in Kaplan et al.'s setup
to 0.50 (RefinedWeb). Counting the output head's parameters, a correction of input measurement, accounts for 39--47
percent of the fall; re-tuning warmup, learning rate, schedule and batch size, flexible inputs that were not
re-optimized at each scale, accounts for the remaining 53--61 percent (Online Appendix~\ref{app:additional}). Both problems are familiar
from production-function estimation \citep{collardwexler2016production,ackerberg2015identification}.

Data quality moves allocation through one channel. Better data can lower the irreducible loss, scale reducible loss
neutrally, or tilt the ratio $A/B$; only the tilt changes the isoquant map, $M^*$ and the wedge. Across Gadre et al.'s
three web corpora, all evaluated on C4, corpus-specific asymptotes with a neutral scale are not rejected ($p=0.90$),
an asymptote shift alone is ($p=0.004$), and the implied tilts move $M^*$ by only 8--18 percent, although the tests
have little power. DataDecide's 25 recipes have the power, and every restriction is rejected, yet the tilt adds only
0.05 points to the 99.7 percent of cross-recipe variation in fit that asymptotes and a neutral scale explain. The tilt
is small in loss but not in allocation: it spans 0.22--0.26 across recipes, which moves $M^*$ 2.9--3.4-fold and the
wedge at given inputs 1.25--1.29-fold, and it is carried by the heavily filtered DCLM recipes. DataDecide identifies the
tilt from checkpoints taken before the schedule ends, and on a common validation set part of any corpus difference is
distribution match (Online Appendix~\ref{app:add-tech}). DeepSeek's IsoFLOP experiments point the same way: their
allocation exponent rose from 0.45 to 0.58 as data quality improved \citep{bi2024deepseek}.

The allocation exponent and $M^*$ are thus properties of a recipe, a tokenizer and a counting convention as much as of
the technology. Section~\ref{sec:wedge} uses developers' own technologies where they exist and reports the range across
all technologies; Section~\ref{sec:econ} shows that $a$, not $\sigma^*$, drives forecasts of data demand.

\subsection{A Controlled Experiment}\label{sec:tech:experiment}

The public designs confound the technology with recipes, units and schedules, and none offers annealed runs at the
tokens-per-parameter ratios of released small models. We therefore ran our own (Online Appendix~\ref{app:sweep}). Two
corpora, FineWeb-Edu and FineWeb (sample-10BT subsets; \citealp{penedo2024fineweb}), draw on the same English web crawl
and differ in filtering: FineWeb-Edu keeps the documents that an educational-quality classifier rates highly. Both share
one 8,192-token byte-level BPE tokenizer, so losses are in common units. The main grid crosses eight decoder-only
transformers of width 128 to 640, with 0.39M--49M non-embedding parameters (1.4M--54M with the tied embedding), with
token budgets from 25M to 800M: 44 endpoints per corpus. $M$ runs from 0.51 to 2,031 with non-embedding parameters and
from 0.46 to 555 with total parameters, and the off-path spread, 1.96, matches Farseer's. Each width has one
constant-learning-rate trunk, from which a copy branches off for each budget and cools down to zero
\citep{hagele2024scaling}, so every endpoint is an annealed run. Training uses contexts of 256 tokens, batches of 16,384
tokens and bf16 matrix products on one Apple M5 Max \citep{hannun2023mlx}. Four-layer models trained to 3.2B tokens
raise the largest $M$ to 4,063 (1,743 with total parameters), so that high $M$ is not reached only by a two-layer
network; learning-rate checks at the corners of the grid, seed replicates in both corpora and evaluation on
WikiText-103 \citep{merity2016pointer}, drawn from neither corpus, address the other threats specific to small
experiments. The largest run uses about $10^{17}$ FLOP, five orders of magnitude below Chinchilla's largest, so the
experiment tests functional form, neutrality and the direction of extrapolation errors, not production-scale levels.

We fixed the analysis before looking at the results. The pre-analysis plan was committed to the project's
version-control repository on 24 September 2026 \citep[on such plans, see][]{olken2015promises}, when part of the
FineWeb-Edu grid had finished but no model had been fitted and no statistic computed. The primary output is loss in bits
per byte, and a difference between corpora counts as quality only if it has the same sign on the other corpus's
validation set and on WikiText-103. Every result is reported with both parameter counts. The estimators are the
model-free $\sigma^*$ and local wedge of subsection~\ref{sec:tech:modelfree} and Huber fits of the Chinchilla form and
the $\kappa$ family; a conclusion is called robust only if it holds both under a wild cluster bootstrap by width, the
unit that shares a trunk \citep{webb2023reworking}, and under a parametric bootstrap with the within-trunk covariance
estimated from seed replicates. The plan poses six questions, each with a decision rule. Data quality is
``factor-biased'' if the 95 percent interval of the tilt excludes zero with the same sign on both validation sets,
``neutral'' if it contains zero on both with a half-width below 0.10, and inconclusive otherwise. The extrapolation
test fits each parametric form on $M\le100$, predicts the wedge at every endpoint with $M>100$, and reports the slope of
$\ln(w_{\rm restricted}/w_{\rm local})$ in $\ln M$, whose sign is the object. One deviation, made for time after the
plan was committed, reduces the reach of the extrapolation test on FineWeb and changes no estimator or decision rule
(Online Appendix~\ref{app:sweep}). Figure~\ref{fig:experiment} summarizes the results.

% TBD-m9: the six questions and their statistics are specified in paper/notes/m9_preanalysis_plan.md.
\textcolor{red}{[TBD-m9: completed endpoints; the power calculation, dated before the results. Q1: model-free and
$\kappa$-free $\sigma^*$ by corpus, both parameter counts, against 0.70 and Porian et al.'s 0.51. Q2: tilt between
corpora on the three validation sets, with the decision and the implied factors for $M^*$ and $\hat w$. Q3: sign of
the extrapolation slope, against Farseer (Section~\ref{sec:wedge}). Q4--Q6: seed noise, learning-rate excess loss
and the bound on $a$, on-path versus full-grid profile likelihood.]}

\begin{figure}[tp]
\centering
\fbox{\parbox{0.9\textwidth}{\centering\vspace{2em}\textcolor{red}{[TBD-m9: two panels. (a) Loss in bits per byte
against $D$ by width for both corpora on a common validation set, with fitted curves and the four-layer high-$M$
runs. (b) Model-free $\sigma^*$ by budget and the local wedge against $\ln M$ for both corpora, with the $M\le100$
extrapolations. File: output/figures/m9\_*.]}\vspace{2em}}}
\caption{A Controlled Experiment: Two Corpora That Differ in Data Quality}\label{fig:experiment}
\begin{figurenotes}
\textcolor{red}{[TBD-m9: design, output, parameter convention, estimator and bootstrap, as pre-registered.]}
\end{figurenotes}
\end{figure}
